\documentclass[10pt]{article}
\usepackage[margin=2cm,landscape=false]{geometry}
\usepackage[T1]{fontenc}
\usepackage{lmodern}
\usepackage{longtable}
\usepackage{array}
\usepackage{xcolor}
\usepackage[hidelinks]{hyperref}
\setlength{\parindent}{0pt}
\setlength{\parskip}{4pt}
\begin{document}
\title{Data model: chinook}
\date{}
\maketitle
\small

Exported by relsynth export-model (format 1); fit seed 7; 3 posterior draws plus the posterior mean in model.json.

\section*{Verification}

7 of 7 checks pass; every condition of the propositions holds.

FAIL: a condition under which a proposition holds is violated.  WARN: a convergence diagnostic of an MCMC block (the model is valid; a longer chain gives more certain draws).

{\footnotesize\begin{longtable}{p{3.2cm}p{1cm}p{7.2cm}p{4.2cm}}
check & result & detail & proved by \\ \hline
parameter supports (mean and 3 draws) & ok & every probability vector is on the simplex, every probability in [0, 1], every rate $\geq$ 0, every quantile function non-decreasing & proper conjugate posteriors: Props 'Block-rate posterior', 'Spike-and-slab Bernoulli', 'Sparse Dirichlet, exact inclusion', 'Posterior mean of a segment marginal' \\
copula correlation matrices (4 parameter sets) & ok & symmetric, unit diagonal, positive semi-definite; smallest eigenvalue 5.28e-03 & Sklar's theorem; Props 'Marginal preservation', 'The sampler targets the copula posterior' \\
tree Employee: size and depth & ok & mean-matrix spectral radius 0.578 (subcritical); the forest is generated conditioned on its 3 level x segment counts, so its size and depth are those of the source (scaled) whatever the radius & Props 'Generation sizes, finiteness and expected size', 'Hierarchy of a given size', 'Hierarchies generated level by level' \\
InvoiceLine: scale rule & ok & every block with a positive rate has source partner rows (m\_l \textgreater{} 0); 0 blocks without & Prop. 'Scale rule' (eta* = eta m\_l / m*\_l keeps rows per anchor under resizing) \\
PlaylistTrack: scale rule & ok & every block with a positive rate has source partner rows (m\_l \textgreater{} 0); 0 blocks without & Prop. 'Scale rule' (eta* = eta m\_l / m*\_l keeps rows per anchor under resizing) \\
latent-segment Gibbs sampler: stationarity & ok & last 15 of 30 final sweeps (after the merge and a burn-in of half the final sweeps; the partition is the last one): split R-hat of the attribute log-likelihood 0.936, Geweke z n/a (stationary: R-hat $\leq$ 1.1, abs(z) \textless{} 3) & blocked Gibbs sampler (layer 1 algorithms paper, Algorithm 3); Geweke (1992); Gelman et al. (2013) \\
copula MCMC: convergence & ok & 20 of 20 chains pass: every correlation has split R-hat $\leq$ 1.1 or a Monte Carlo standard error $\leq$ 0.01 on the Fisher-z scale (largest R-hat 1.096, largest standard error 0.1684) & Prop. 'Rank likelihood with parameter expansion'; split R-hat (Gelman et al. 2013); fixed-width stopping (Jones et al. 2006; Flegal, Haran \& Jones 2008) \\
\end{longtable}}

\textbf{Replica against the source} (scorecard of \texttt{relsynth} on the replica's output, posterior mean, same checks as \texttt{relsynth run}):

\begin{itemize}
\item replica scorecard: 139/140 checks pass
\item replica rows: Artist 275, Employee 8, Genre 25, MediaType 5, Playlist 18, Album 347, Customer 59, Invoice 412, Track 3,503, InvoiceLine 2,178, PlaylistTrack 8,710
\item seconds: 1.5
\end{itemize}

\section*{1. The generative model}

A synthetic database D* is a draw from the posterior predictive distribution

\begin{verbatim}
theta ~ p(theta | D, z),   D* ~ p(D* | theta)
\end{verbatim}
where z are the latent segments fitted by the blocked Gibbs sampler (one segment per
row of every entity table) and theta collects every parameter block below.  All
blocks are conjugate given z, so p(theta | D, z) is a product of the closed-form
posteriors listed here; draw m of model.json is the theta that \texttt{relsynth run} uses
for synthetic dataset m.  Tables are generated parents first (section 3).

\textbf{Entity table e} (n\_e rows, K\_e segments)

\begin{verbatim}
pi_e ~ Dirichlet(1 / K_e + counts of rows per segment)               (segment shares)
segment sizes: n*_e pi_e rounded by largest remainder, n*_e = scale x n_e
row i of segment k: its modelled columns x_i ~ copula_k (below)
\end{verbatim}
A self-referencing table (a parent-child tree) is a multi-type Galton-Watson
process typed by (segment, depth):

\begin{verbatim}
lam_{k,d} ~ Gamma(c lam_k + children of depth-d nodes of segment k, c + those nodes)
psi_{k,d} ~ Dirichlet(c psi_k + child-segment counts)
the forest is drawn with round(scale x n_{k,d}) nodes of segment k at depth d; a
node of depth d picks its parent at depth d-1 with probability
proportional to w_i lam_{k_i,d-1} psi_{k_i,d-1}(k)   (w_i ~ Gamma(alpha_k, alpha_k) frailty)
\end{verbatim}
\textbf{Columns of a table given a group g} (a segment for entities, a segment block
for facts): a Gaussian copula with correlation matrix R\_g over the latent
normal scores of the modelled columns, and one marginal per column:

\begin{verbatim}
R_g ~ p(R | ranks)  (LKJ(eta) prior, rank likelihood; MCMC draws kept)
discrete column:   p_g ~ Dirichlet(a p_coarse + level counts of the group)
continuous column: F_g ~ DP(a + n_g, (a F_coarse + sum_i delta_{x_i}) / (a + n_g))
                   (Bayesian bootstrap of the group mixed with its coarse group's law;
                    stored as 49 quantile knots)
null share:        discrete q_g ~ Beta(a q_coarse + nulls, a (1 - q_coarse) + non-nulls),
                   continuous q_g ~ Beta(1 + nulls, 1 + non-nulls)
                   (with point masses at 0 and 1 under the structural prior)
sampling: Z ~ N(0, R_g), U = Phi(Z), x_j = F_j^{-1}(U_j)  (null when U_j < q_j)
\end{verbatim}
Columns joined into one categorical (functional dependencies, strongly
associated pairs) are one discrete column over their observed combinations.

\textbf{Many-to-one link child $\rightarrow$ parent}

\begin{verbatim}
eta_{kl} ~ Gamma(1/2 + links from segment k to segment l, b + n_k m_l),  b = 1 / (2 x overall link rate)
           (with a point mass at 0 under the structural prior)
phi_{k.} = eta_{k.} m_. / sum_l eta_{kl} m_l           (parent segment law of a child of segment k)
p_null,k ~ Beta(nulls of segment k + small prior, non-nulls + 1)
parent row within segment l: by degree weight w_j (Gamma(alpha_l, alpha_l) frailty,
alpha_l by the method of moments when the counts are over-dispersed), or by the
DP posterior predictive of the counts per parent row (Bayesian bootstrap of the
source counts with mass a on Poisson(mean)), conditioned on the segment's total
\end{verbatim}
\textbf{Facts}

\begin{verbatim}
single dimension:  rows of anchor row i ~ Poisson(lam_{z_i} w_i), lam_k ~ Gamma(a + rows, b + anchors)
                   (or 1 + Poisson, or the DP count law)
Poisson block:     N_{kl} ~ Poisson(eta*_{kl} W1_k W2_l), eta_{kl} ~ Gamma(a + C_{kl}, b + E_{kl}),
                   E_{kl} = n_k m_l, eta* = eta m_l / m*_l (scale rule), rows by weight within the blocks
occupancy:         anchor active ~ Bernoulli(pi_k), pair present ~ Bernoulli(p*_{kl}),
                   pi_k ~ Beta(a + active, b + inactive), p_{kl} ~ Beta(a + present, b + absent pairs at risk)
                   (+ point masses at 0 and 1), mu ~ Gamma(a_m + extra rows, b_m + present pairs),
                   p* = min(1, p m_l / m*_l); rows per present pair 1 + Poisson(mu)
further keys:      segment of the dimension row ~ P(. | block) ~ sparse Dirichlet
                   backed off to the more informative block dimension
time column:       period ~ theta_{z_i}, theta_k ~ Dirichlet(period counts + 2 x global profile),
                   uniform within the period; null with q ~ Beta(1 + nulls, 1 + rest)
columns:           the block's copula and marginals; blocks with few rows use the
                   coarse group of the more informative dimension; parent attributes
                   (condition_on_parents) enter as observed copula coordinates
                   (conditional normal)
\end{verbatim}
\textbf{Computed columns}: lookups (a parent's column through a key), derived columns
(expressions) and aggregates (a parent column computed from its children) are
deterministic functions of the generated columns, evaluated in dependency order.

\section*{2. This dataset's model (posterior means)}

\subsection*{Artist (entity, 275 rows)}
Segments: K = 2; shares 0.955, 0.045


\subsection*{Employee (entity, 8 rows)}
Segments: K = 5; shares 0.133, 0.133, 0.356, 0.133, 0.244

Strata: segments nested in Title (5 strata); a row's stratum columns follow from its segment.

Tree Employee.ReportsTo$\rightarrow$Employee: 3 levels, nodes per level 1, 2, 5.

Columns (3 modelled; table-level marginals):


\begin{itemize}
\item BirthDate: continuous (date); median 1966-08-06, 5\%-95\% 1947-09-19 to 1973-08-29
\item HireDate: continuous (date); median 2003-07-25, 5\%-95\% 2002-04-01 to 2004-03-04
\item City: discrete, 3 levels; Calgary: 0.616, Lethbridge: 0.253, Edmonton: 0.131
\end{itemize}
Strongest latent correlations (pooled): BirthDate \textasciitilde{} HireDate -0.22, BirthDate \textasciitilde{} City +0.00, HireDate \textasciitilde{} City +0.00

Groups with their own marginals: 5; own copulas: 0.



\subsection*{Genre (entity, 25 rows)}
Copied reference table: its rows are kept as they are; 25 latent segments (one per row when small) carry its links.


\subsection*{MediaType (entity, 5 rows)}
Copied reference table: its rows are kept as they are; 5 latent segments (one per row when small) carry its links.


\subsection*{Playlist (entity, 18 rows)}
Copied reference table: its rows are kept as they are; 18 latent segments (one per row when small) carry its links.


\subsection*{Album (entity, 347 rows)}
Segments: K = 8; shares 0.259, 0.069, 0.164, 0.107, 0.098, 0.035, 0.066, 0.202

Link ArtistId $\rightarrow$ Artist: parent segment law phi (degree weights); null share by segment 0.000, 0.000, 0.000, 0.000, 0.000, 0.000, 0.000, 0.000


{\footnotesize\begin{longtable}{lll}
 & 0 & 1 \\ \hline
0 & 1.000 & 0.000 \\
1 & 1.000 & 0.000 \\
2 & 0.394 & 0.606 \\
3 & 0.652 & 0.348 \\
4 & 0.967 & 0.033 \\
5 & 0.935 & 0.065 \\
6 & 1.000 & 0.000 \\
7 & 0.536 & 0.464 \\
\end{longtable}}


\subsection*{Customer (entity, 59 rows)}
Segments: K = 1; shares 1.000

Columns (1 modelled; table-level marginals):


\begin{itemize}
\item City|State|Country: discrete, 53 levels; (Prague, null, Czech Republic): 0.034, (São Paulo, SP, Brazil): 0.034, (Mountain View, CA, USA): 0.034, (Berlin, null, Germany): 0.034
\end{itemize}
Groups with their own marginals: 1; own copulas: 0.


Link SupportRepId $\rightarrow$ Employee: parent segment law phi (degree weights); null share by segment 0.000


{\footnotesize\begin{longtable}{llllll}
 & 0 & 1 & 2 & 3 & 4 \\ \hline
0 & 0.000 & 0.000 & 1.000 & 0.000 & 0.000 \\
\end{longtable}}


\subsection*{Invoice (entity, 412 rows)}
Segments: K = 8; shares 0.133, 0.148, 0.126, 0.131, 0.119, 0.245, 0.032, 0.066

Columns (1 modelled; table-level marginals):


\begin{itemize}
\item InvoiceDate: continuous (date); median 2011-06-23, 5\%-95\% 2009-03-26 to 2013-09-27
\end{itemize}
Groups with their own marginals: 8; own copulas: 0.


Link CustomerId $\rightarrow$ Customer: parent segment law phi (DP counts given the total); null share by segment 0.000, 0.000, 0.000, 0.000, 0.000, 0.000, 0.000, 0.000


{\footnotesize\begin{longtable}{ll}
 & 0 \\ \hline
0 & 1.000 \\
1 & 1.000 \\
2 & 1.000 \\
3 & 1.000 \\
4 & 1.000 \\
5 & 1.000 \\
6 & 1.000 \\
7 & 1.000 \\
\end{longtable}}


\subsection*{Track (entity, 3,503 rows)}
Segments: K = 52; shares 0.046, 0.077, 0.056, 0.080, 0.053, 0.025, 0.009, 0.024, 0.007, 0.029, 0.047, 0.060, 0.057, 0.038, 0.003, 0.009, 0.014, 0.052, 0.052, 0.043 ...

Strata: segments nested in MediaTypeId, GenreId (38 strata); a row's stratum columns follow from its segment.

Columns (3 modelled; table-level marginals):


\begin{itemize}
\item Milliseconds: continuous (integer); median 2.556e+05, 5\%-95\% 1.419e+05 to 1.571e+06
\item Bytes: continuous (integer); median 8.108e+06, 5\%-95\% 3.742e+06 to 2.103e+08
\item UnitPrice: discrete, 2 levels; 0.99: 0.939, 1.99: 0.061
\end{itemize}
Strongest latent correlations (pooled): Milliseconds \textasciitilde{} Bytes +0.88, Bytes \textasciitilde{} UnitPrice +0.02, Milliseconds \textasciitilde{} UnitPrice +0.01

Groups with their own marginals: 52; own copulas: 16.


Link AlbumId $\rightarrow$ Album: parent segment law phi (every parent has a child (1 + Poisson)); null share by segment 0.000, 0.000, 0.000, 0.000, 0.000, 0.000, 0.000, 0.000, 0.000, 0.000, 0.000, 0.000


52 x 8 matrix; 416 positive entries, range 1.2e-299 to 1



\subsection*{InvoiceLine (fact, poisson model, anchor Invoice, partner Track)}
Block rates eta (anchor segment x partner segment):


8 x 52 matrix; 416 positive entries, range 1e-300 to 0.0245


Conditioned on parent columns: Track.Milliseconds, Track.Bytes, Invoice.InvoiceDate.

Columns (4 modelled; table-level marginals):


\begin{itemize}
\item Quantity: discrete, 1 levels; 1: 1.000
\item \_\_aux\_\_TrackId.Milliseconds (parent column, conditioning only): continuous (integer); median 2.564e+05, 5\%-95\% 1.427e+05 to 1.136e+06
\item \_\_aux\_\_TrackId.Bytes (parent column, conditioning only): continuous (integer); median 8.105e+06, 5\%-95\% 3.891e+06 to 1.724e+08
\item \_\_aux\_\_InvoiceId.InvoiceDate (parent column, conditioning only): continuous (date); median 2011-06-24, 5\%-95\% 2009-03-20 to 2013-09-29
\end{itemize}
Strongest latent correlations (pooled): \_\_aux\_\_TrackId.Milliseconds \textasciitilde{} \_\_aux\_\_TrackId.Bytes +0.87, \_\_aux\_\_TrackId.Milliseconds \textasciitilde{} \_\_aux\_\_InvoiceId.InvoiceDate +0.00, \_\_aux\_\_TrackId.Bytes \textasciitilde{} \_\_aux\_\_InvoiceId.InvoiceDate -0.00, Quantity \textasciitilde{} \_\_aux\_\_TrackId.Milliseconds +0.00, Quantity \textasciitilde{} \_\_aux\_\_InvoiceId.InvoiceDate +0.00

Groups with their own marginals: 71; own copulas: 0.



\subsection*{PlaylistTrack (fact, occupancy model, anchor Playlist, partner Track)}
Anchor activity pi by segment: 1.000, 0.000, 1.000, 0.000, 1.000, 0.000, 0.000, 1.000, 1.000, 1.000, 1.000, 1.000, 1.000, 1.000, 1.000, 1.000, 1.000, 1.000

Presence probabilities p (anchor segment x partner segment); multiplicity by none:


18 x 52 matrix; 368 positive entries, range 0.00684 to 1


Conditioned on parent columns: Track.Milliseconds, Track.Bytes.

Columns (2 modelled; table-level marginals):


\begin{itemize}
\item \_\_aux\_\_TrackId.Milliseconds (parent column, conditioning only): continuous (integer); median 2.56e+05, 5\%-95\% 1.398e+05 to 1.135e+06
\item \_\_aux\_\_TrackId.Bytes (parent column, conditioning only): continuous (integer); median 8.102e+06, 5\%-95\% 3.59e+06 to 1.701e+08
\end{itemize}
Strongest latent correlations (pooled): \_\_aux\_\_TrackId.Milliseconds \textasciitilde{} \_\_aux\_\_TrackId.Bytes +0.88

Groups with their own marginals: 83; own copulas: 0.



\section*{3. Generation order}

Artist $\rightarrow$ Employee $\rightarrow$ Genre $\rightarrow$ MediaType $\rightarrow$ Playlist $\rightarrow$ Album $\rightarrow$ Customer $\rightarrow$ Invoice $\rightarrow$ Track $\rightarrow$ InvoiceLine $\rightarrow$ PlaylistTrack

\section*{4. Computed columns}

\begin{itemize}
\item Invoice.BillingAddress = CustomerId.Address
\item Invoice.BillingCity = CustomerId.City
\item Invoice.BillingState = CustomerId.State
\item Invoice.BillingCountry = CustomerId.Country
\item Invoice.BillingPostalCode = CustomerId.PostalCode
\item InvoiceLine.UnitPrice = TrackId.UnitPrice
\item Invoice.Total = sum(UnitPrice * Quantity) over InvoiceLine by InvoiceId
\end{itemize}

\section*{5. Replicating the synthesis}

\begin{itemize}
\item Exactly, with relsynth: \texttt{python run\_model.py} (the calls of \texttt{relsynth run}; same spec, same seed, bit-identical output).
\item On your own, without relsynth: \texttt{python replicate.py model.json --out replica/ [--draw 0|mean] [--seed 7] [--scale 1.0] [--reference DIR]} (numpy, pandas, scipy).  It draws the core process of section 1 from the exported posterior draws.
\end{itemize}

What replicate.py leaves to the engine:

\begin{itemize}
\item Album: Title (text, identifiers and sequences; the engine's text and identifier models are not exported because they encode source strings)
\item Artist: Name (text, identifiers and sequences; the engine's text and identifier models are not exported because they encode source strings)
\item Customer: FirstName, LastName, Company, Address, PostalCode, Phone, Fax, Email (text, identifiers and sequences; the engine's text and identifier models are not exported because they encode source strings)
\item Employee: LastName, FirstName, Address, PostalCode, Phone, Fax, Email (text, identifiers and sequences; the engine's text and identifier models are not exported because they encode source strings)
\item Track: Name, Composer (text, identifiers and sequences; the engine's text and identifier models are not exported because they encode source strings)
\item Invoice: keys numbered in the order of InvoiceDate (engine only)
\end{itemize}
\end{document}
